\chapter{Descrizione quantistica dello Stern-Gerlach}

Adesso che abbiamo introdotto il formalismo quantistico per descrivere sistemi tramite le matrici densità, possiamo andare a dare una descrizione quantistica di quanto osservato negli esperimenti dello stern-gerlach. 

L'atomo sparato dal forno nell'apparato ha due tipi di gradi di libertà: dove si trova, cioè le varibili orbitali, posizione e impulso e il suo grado di libertà interno, lo spin. \\
Queste due varibili sono cinematicamente indipendenti, quindi posso descrivere lo spazio di Hilbert del sistema come 
\begin{tcolorbox}[colback=black!5!white, colframe=black!45!white, boxrule=0.8pt]
\begin{equation}
    \mathcal{H} = \mathcal{H}_{\text{int}} \otimes \mathcal{H}_{\text{orb}}
\end{equation}
\end{tcolorbox}
In particolare di ha che $\text{dim} \ \mathcal{H}_{\text{orb}} = \infty$, in quanto la posizione è continua, e $\text{dim} \ \mathcal{H}_{\text{int}} = 2$, in quanto lo spin può essere solo \emph{up} o \emph{down}. 

\section{Spazio $\mathcal{H}_{\text{int}}$}

Vogliamo ora descrivere esplicitamente lo spazio $\mathcal{H}_{\text{int}}$, individuandone una base e ricavando le relazioni fra le basi associate ad assi diversi.

Indichiamo con $\ket{z, +}$ e $\ket{z, -}$ gli stati delle particelle cheemergono rispettivamente dal canale superiore e da quello inferiore di un apparato orientato lungo $\hat{z}$.

Questi due stati sono ortogonali, e il motivo è sperimentale: una sorgente pura $\ket{z, +}$ rianalizzata lungo $\hat{z}$ restituisce il $100\%$ degli atomi nel canale up e nessuno in quello down. Probabilità nulla significa prodotto scalare nullo,
%
\begin{equation}
    \braket{z, + |z, -} = 0
    \label{eq:ortogonalita-z}
\end{equation}
%
Essendo ortogonali sono in particolare linearmente indipendenti, e poiché $\dim \mathcal{H}_{\text{int}} = 2$ essi costituiscono una base ortonormale:
%
\begin{equation}
    \mathcal{H}_{\text{int}} = \operatorname{span}
    \bigl\{ \ket{z, +} \, , \, \ket{z, -} \bigr\}
    \label{eq:span-int}
\end{equation}

Ogni stato interno è dunque rappresentato da un vettore della forma
%
\begin{equation}
    \ket{\psi} = \alpha \ket{z, +} + \beta \ket{z, -}
    \, , \qquad |\alpha|^2 + |\beta|^2 = 1
    \label{eq:stato-interno}
\end{equation}
%
Conviene però riscriverlo in una forma più maneggevole. Poiché lo stato è definito a meno di una fase globale, possiamo sfruttare questa libertà per rendere reale e positivo il primo coefficiente, ottenendo
%
\begin{equation}
    \ket{\psi} = |\alpha| \ket{z, +}
    + \sqrt{1 - |\alpha|^2} \; e^{i \theta} \ket{z, -}
    \label{eq:stato-interno-ridotto}
\end{equation}
%
dove restano due soli parametri reali: il modulo $|\alpha|$ e la fase relativa $\theta$.

\subsection*{Stati lungo $\hat{x}$ e $\hat{y}$}

Disponendo della base lungo $\hat{z}$, possiamo ricavare la forma esplicita degli stati associati agli altri assi.

Il dato sperimentale di partenza è che una sorgente $\ket{z, +}$ inviata in un filtro orientato lungo $\hat{x}$ si ripartisce equamente fra i due canali, e altrettanto accade per una sorgente $\ket{z, -}$:
%
\begin{equation*}
    \bigl| \braket{x, + |z, +} \bigr|^2 = \bigl| \braket{x, - |z, +} \bigr|^2 = \frac{1}{2}
\end{equation*}
%
Questa informazione fissa il \emph{modulo} dei coefficienti dello sviluppo, ma non la loro fase, che resta per il momento indeterminata. Scriviamodunque
%
\begin{equation*}
    \ket{z,+} = \frac{\ket{x,+}}{\sqrt{2}}
    + e^{i\theta_+} \frac{\ket{x,-}}{\sqrt{2}}
    \, , \qquad
    \ket{z,-} = \frac{\ket{x,+}}{\sqrt{2}}
    + e^{i\theta_-} \frac{\ket{x,-}}{\sqrt{2}}
\end{equation*}

Le due fasi non sono indipendenti: i vettori $\ket{z,+}$ e $\ket{z,-}$ devono continuare a soddisfare la \eqref{eq:ortogonalita-z}. Imponendola,
%
\begin{equation*}
    \braket{z,-|z,+} = \frac{1}{2} + \frac{1}{2} e^{i(\theta_+ - \theta_-)} = 0 \implies e^{i(\theta_+ - \theta_-)} = -1 \implies \theta_+ - \theta_- = \pi
\end{equation*}

Resta determinata soltanto la \emph{differenza} delle due fasi: quella comune corrisponde a ridefinire la fase del vettore di base $\ket{x,-}$, ed è dunque una convenzione priva di contenuto fisico. Fissiamola scegliendo $\theta_- = 0$ e $\theta_+ = \pi$, cosicché
%
\begin{equation}
\label{eq:relazioni_z_x}
    \ket{z,+} = \frac{\ket{x,+} - \ket{x,-}}{\sqrt{2}} \quad , \quad \ket{z,-} = \frac{\ket{x,+} + \ket{x,-}}{\sqrt{2}}
\end{equation}
%
e invertendo queste due relazioni si ottiene infine
%
\begin{tcolorbox}[colback=black!5!white, colframe=black!45!white, boxrule=0.8pt]
\begin{equation}
    \label{eq:base-x}
    \ket{x, +} = \frac{\ket{z, +} + \ket{z, -}}{\sqrt{2}}\quad , \quad \ket{x, -} = \frac{\ket{z, -} - \ket{z, +}}{\sqrt{2}}
\end{equation}
\end{tcolorbox}

Possiamo fare un ragionamente simile anche per un filtro orientato lungo $\hat{y}$, ottenendo 
\begin{equation*}
    \ket{y, \pm} = \frac{\ket{z,+} + e^{i\phi_{\pm}}\ket{z,-}}{\sqrt{2}} \quad , \quad \phi_+ - \phi_- = \pi
\end{equation*}
Questa volta, però, dobbiamo anche considerare il fatto che, facendo passare una sorgente $\ket{x, \pm}$ in un filtro orientato lungo $\hat{y}$ si ripartisce equamente tra i due canali
\begin{equation*}
    |\braket{y,+|x,+}|^2=\left|\frac{1+e^{-i\phi_+}}{2}\right|^2=\frac{1+\cos\phi_+}{2}=\frac12 \implies \cos\phi_+=0 \implies e^{i\phi_+}=\pm i
\end{equation*}
Prendendo $\phi_+=\pi/2$ e quindi $\phi_-=-\pi/2$:
\begin{tcolorbox}[colback=black!5!white, colframe=black!45!white, boxrule=0.8pt]
\begin{equation}
    \ket{y,+}=\frac{\ket{z,+} + i\,\ket{z,-}}{\sqrt2} \quad , \quad \ket{y,-}=\frac{\ket{z,+} - i\,\ket{z,-}}{\sqrt2}
\end{equation}
\end{tcolorbox}

\newpage

\section{Ricollimazione del fascio}

Adesso che abbiamo trovato un modo per identificare gli stati della particella, possiamo andare a spiegare i risultati trovati durante l'esperimento, in particolare in questa sezione andremo a trattare la ricollimazione del fascio e vedremo in dettaglio il perché la presenza di un rivelatore modifichi lo stato di uscita.

In questo caso un sistema misto può essere scritto come 
\begin{equation}
    \hat{\rho}_{\text{misto}} = \frac{\ketbra{z, +}{z, +} + \ketbra{z, -}{z, -}}{2} \otimes \ketbra{\psi}{\psi} = \frac{\hat{\mathds{1}}}{2} \otimes \ketbra{\psi}{\psi}
\end{equation}
Oppure, al posto di $z$, mettendo $x$ o $y$.

\subsection*{L'azione dell'apparato}

Tutto quanto segue discende da una sola ipotesi fisica, che riguarda il modo
in cui il campo disomogeneo agisce sull'atomo. Detto $\ket{\text{in}}$ lo
stato orbitale del fascio collimato in ingresso, un apparato orientato lungo
$\hat{x}$ separa spazialmente le due componenti dello spin:
%
\begin{tcolorbox}[colback=black!5!white, colframe=black!45!white, boxrule=0.8pt]
\begin{equation}
    \hat{U}_1 : \quad
    \ket{x,+} \otimes \ket{\text{in}} \longmapsto \ket{x,+} \otimes \ket{S}
    \quad , \quad
    \ket{x,-} \otimes \ket{\text{in}} \longmapsto \ket{x,-} \otimes \ket{I}
    \label{eq:apparato}
\end{equation}
\end{tcolorbox}
%
dove $\ket{S}$ e $\ket{I}$ sono gli stati orbitali dei due cammini, superiore
e inferiore, e sono ortogonali, $\braket{S|I} = 0$, perché corrispondono a
regioni di spazio distinte.

Si noti che la \eqref{eq:apparato} definisce $\hat{U}_1$ sui soli vettori
della base di $\hat{x}$: su una generica sovrapposizione agisce per
\emph{linearità}. È da questa semplice circostanza che discenderà tutto il
resto.

Il secondo apparato, di gradiente invertito, riporta i due cammini a
sovrapporsi in un unico stato orbitale $\ket{U}$:
%
\begin{tcolorbox}[colback=black!5!white, colframe=black!45!white, boxrule=0.8pt]
\begin{equation}
    \hat{U}_2 : \quad
    \ket{x,+} \otimes \ket{S} \longmapsto \ket{x,+} \otimes \ket{\text{out}}
    \quad , \quad
    \ket{x,-} \otimes \ket{I} \longmapsto \ket{x,-} \otimes \ket{\text{out}}
    \label{eq:ricombinatore}
\end{equation}
\end{tcolorbox}
%
Anche questa è una trasformazione unitaria: i due vettori di partenza sono
ortogonali e tali restano quelli di arrivo, che differiscono nel fattore di
spin\footnote{Basta che $\hat{U}_2$ sia definita sul sottospazio
effettivamente occupato dopo il primo apparato, cioè
$\operatorname{span}\{\ket{x,+}\ket{S}, \ket{x,-}\ket{I}\}$.}.

\newpage

\subsection{Ricollimazione senza rivelatore}

Consideriamo come prima cosa il caso di una ricollimazione quando \emph{non} è presente un rivelatore. 

\subsection*{Sorgente mista}

\begin{figure}[htbp]
\centering
\begin{tikzpicture}[sgstyle]
  \ricombinazione{$\hat{\rho}_{\text{misto}}$}
    {$\ketbra{x,+}{x,+}$}{$\ketbra{x,-}{x,-}$}
    {$\dfrac{\hat{\mathds{1}}}{2}$}
\end{tikzpicture}
\caption{Ricombinazione con sorgente mista: lo stato in uscita coincide
con quello in ingresso.}
\label{fig:ricomb-misto}
\end{figure}

Lo stato iniziale è
\begin{equation*}
    \hat{\rho}_{\text{in}} = \frac{\hat{\mathds{1}}}{2} \otimes \ketbra{\text{in}}{\text{in}}
    = \left( \frac{\ketbra{x,+}{x,+} + \ketbra{x,-}{x,-}}{2} \right)
    \otimes \ketbra{\text{in}}{\text{in}}
\end{equation*}
dove nella seconda uguaglianza si è usato che $\hat{\mathds{1}}/2$ è la stessa matrice in ogni base, e in particolare in quella del filtro attraversato.

Applicando la \eqref{eq:apparato} a ciascuno dei due addendi si ottiene lo stato fra i due apparati:
\begin{equation}
    \hat{\rho}_1 = \frac12 \, \ketbra{x,+}{x,+} \otimes \ketbra{S}{S}
    + \frac12 \, \ketbra{x,-}{x,-} \otimes \ketbra{I}{I}
    \label{eq:misto-intermedio}
\end{equation}
Applicando infine la \eqref{eq:ricombinatore}, entrambi i fattori orbitali diventano $\ketbra{\text{out}}{\text{out}}$, che si raccoglie:
\begin{equation}
    \hat{\rho}_{\text{out}} = \frac{\hat{\mathds{1}}}{2} \otimes \ketbra{\text{out}}{\text{out}}
    \label{eq:misto-out}
\end{equation}

Pertanto si ha, come ci aspettavamo, che lo stato di uscita è uno stato misto. 

\subsection*{Sorgente pura}

\begin{figure}[htbp]
\centering
\begin{tikzpicture}[sgstyle]
    \ricombinazione{$\ket{z,+}$}
    {$\tfrac{1}{\sqrt2}\ket{x,+}\ket{S}$}{$-\tfrac{1}{\sqrt2}\ket{x,-}\ket{I}$}
    {$\ket{z,+}$}
\end{tikzpicture}
\caption{Ricombinazione con sorgente pura: lo stato $\ket{z,+}$ viene
ricostruito all'uscita.}
\label{fig:ricomb-puro}
\end{figure}

Lo stato iniziale è 
\begin{equation}
    \hat{\rho}_{\text{in}} = \ketbra{z, +}{z, +} \otimes \ketbra{\text{in}}{\text{in}}
\end{equation}
Sfruttando la relazione \eqref{eq:relazioni_z_x} possiamo scrivere la matrice densità in termini di $\ket{x, \pm}$
\begin{align*}
    \ketbra{z, +}{z,+} &= \frac{(\ket{x, +} - \ket{x, -})(\bra{x, +} - \bra{x, -})}{2}\\
    &= \frac{\ketbra{x,+}{x,+} + \ketbra{x,-}{x, -} - \ketbra{x, +}{x, -} - \ketbra{x, -}{x, +}}{2}
\end{align*}
In questo modo, dopo il primo stern-gerlach, lo stato diventa
\begin{align*}
    \hat{\rho}_1 = \frac{\ketbra{x, +}{x, +} \otimes \ketbra{S}{S} - \ketbra{x,-}{x, +} \otimes \ketbra{I}{S} + \ketbra{x, -}{x, -} \otimes \ketbra{I}{I} - \ketbra{x,+}{x, -} \otimes \ketbra{S}{I}}{2}
\end{align*}
Uscendo, possiamo applicare la \eqref{eq:ricombinatore}, tutti i fattori orbitali diventano $\ketbra{\text{out}}{\text{out}}$, che si raccoglie:
%
\begin{equation}
    \hat{\rho}_{\text{out}} = \ketbra{z, +}{z, +} \otimes \ketbra{\text{out}}{\text{out}}
\end{equation}

Quindi si ha che lo stato di uscita è uno stato puro, come si vede nell'esperimento. In questo modo è conservata la coerenza quantomeccanica. 

\subsection{Ricollimazione con rivelatore}

Consideriamo adesso il caso in cui sul ramo up ci sia un rivelatore, modellizzato come uno stumento a due stati $\ket{0}$ e $\ket{1}$ che in quiete si trova nello stato $\ket{0}$ e quando rileva il passaggio di una particella passa allo stato $\ket{1}$. 

Andiamo ad analizzare il caso in cui la sorgente di partenza è una sorgente pura. Chiamiamo il sistema della sorgente sistema $A$, mentre chiamiamo il sistema del rivelatore sistema $B$. 

\begin{figure}[htbp]
\centering
\begin{tikzpicture}[sgstyle]
  \ricombinazionerivelatore{$\ket{z,+}_A$}
    {$\ket{x,+}\ket{S}$}{$\ket{x,-}\ket{I}$}
    {$\dfrac{\hat{\mathds{1}}}{2}$}
\end{tikzpicture}
\caption{Ricombinazione con un rivelatore $B$ sul ramo superiore. Il
rivelatore passa da $\ket{0}_B$ a $\ket{1}_B$ al passaggio della
particella; lo stato in uscita è misto, pur essendo pura la sorgente.}
\label{fig:ricomb-rivelatore}
\end{figure}

Lo stato iniziale è 
\begin{align*}
    \hat{\rho}_{\text{in}} &= \ketbra{z,+}{z,+}_A \otimes \ketbra{\text{in}}{\text{in}}_A \otimes \ketbra{0}{0}_B\\
    &= \frac{\ketbra{x,+}{x,+}_A + \ketbra{x,-}{x, -}_A - \ketbra{x, +}{x, -}_A - \ketbra{x, -}{x, +}_A}{2} \otimes \ketbra{\text{in}}{\text{in}}_A \otimes \ketbra{0}{0}_B
\end{align*}

Dopo il passaggio per il primo stern-gerlach ma \emph{prima} del passaggio attraverso il rivelatore lo stato è 
\begin{align*}
    \hat{\rho}_{1,i} = \frac{\ketbra{x, +}{x, +} \otimes \ketbra{S}{S} - \ketbra{x,-}{x, +} \otimes \ketbra{I}{S} + \ketbra{x, -}{x, -} \otimes \ketbra{I}{I} - \ketbra{x,+}{x, -} \otimes \ketbra{S}{I}}{2} \otimes \ketbra{0}{0}_B
\end{align*}
Invece, \emph{dopo} il passaggio
\begin{align*}
    \hat{\rho}_{1,f} = \frac{1}{2} \Bigl[ &
      \ketbra{x,+}{x,+} \otimes \ketbra{S}{S} \otimes \ketbra{1}{1}
    + \ketbra{x,-}{x,-} \otimes \ketbra{I}{I} \otimes \ketbra{0}{0} \\
    &- \ketbra{x,+}{x,-} \otimes \ketbra{S}{I} \otimes \ketbra{1}{0}
     - \ketbra{x,-}{x,+} \otimes \ketbra{I}{S} \otimes \ketbra{0}{1}
    \Bigr]
\end{align*}
Infine, all'uscita 
\begin{equation*}
    \hat{\rho}_{\text{out}} = \frac{1}{2} \Bigl[\ketbra{x,+}{x,+} \otimes \ketbra{1}{1} + \ketbra{x,-}{x,-} \otimes \ketbra{0}{0} - \ketbra{x,+}{x,-} \otimes \ketbra{1}{0} - \ketbra{x,-}{x,+} \otimes \ketbra{0}{1}\Bigr] \otimes \ketbra{\text{out}}{\text{out}}
\end{equation*}

Dato che a noi interessa lo stato di $A$, ricaviamo la matrice densità ridotta $\rho_A$:
\begin{equation*}
    \hat{\rho}_A = \operatorname{tr}_B (\hat{\rho}_{\text{out}}) = \frac{\ketbra{x, +}{x, +} + \ketbra{x, -}{x, -}}{2} \otimes \ketbra{\text{out}}{\text{out}} = \frac{\hat{\mathds{1}}}{2} \otimes \ketbra{\text{out}}{\text{out}}
\end{equation*}
Si ha allora in questo caso che pur avendo come sorgente input una sorgente pura, non otteniamo come output uno stato puro, bensì uno stato misto. 